
\section{Introdução}


A análise de comportamento de compra dos consumidores em plataformas de e-commerce tem se tornado uma área central de pesquisa devido à crescente competitividade no mercado digital e à importância de entender as preferências dos clientes para melhorar a experiência de compra e maximizar a rentabilidade. O comportamento do cliente online, especialmente no que diz respeito ao tempo e à frequência com que ele realiza compras, pode fornecer insights valiosos sobre a sua lealdade e intenção de compra. No entanto, prever com precisão quando um cliente fará uma compra ou quando ele se tornará inativo representa um desafio significativo para os profissionais da estatística. 

 O uso de modelos estatísticos ou de modelos treinados a partir de algoritmos de aprendizado de máquina podem ser úteis  para prever o comportamento o que contribui para que as empresas, através de seus gestores, tomem decisões mais acertivas. Isso pode  direcionar suas estratégias de vendas e marketing. A possibilidade de identificar padrões e comportamentos de clientes antes que ocorra uma compra ou um abandono pode ser a chave para otimizar campanhas de marketing e aumentar a retenção de clientes. Além disso, a compreensão dos métodos de análise de sobrevivência, em especial via aprendizado de máquina, para  tratar o tempo de espera até a ocorrência de um evento,como a compra de um produto, de maneira eficaz e incorporando censura e variáveis explicativas contribui grandemente para a formação profissional do estatístico ou analista de dados.

A questão central deste estudo é como modelar o tempo até a próxima compra com base no comportamento ou características do cliente, utilizando dados de um conjunto e-commerce. Além disso, como pode ser previsto  a probabilidade de um cliente não comprar novamente  (abandono) dentro de um intervalo de tempo é outra questão interessante. Esses questionamentos são relevantes pois suas respostas podem contribuir para o entendimento da dinâmica de compra, prever a rotatividade de clientes e otimizar a alocação de recursos de marketing.

 O linguagem de programação Python \cite{python} é uma excelente escolha para implementar essa análise, pois oferece bibliotecas especializadas para análise de sobrevivência, como a \textit{lifelines}, que contém funções tanto para o ajuste de modelos de Cox quanto para o cálculo da curva de Kaplan-Meier. Além disso, o Python é uma plataforma amplamente utilizada em ciência de dados, proporcionando uma grande quantidade de recursos e documentação, o que facilita o aprendizado e a implementação de análises complexas.

 Os dados sobre os quais será aplicada a metodologia vêm do ``E-commerce Customer Behavior and Purchase Datase'', um conjunto de dados sintético gerado pela biblioteca Faker do Python, disponibilizado na plataforma Kaggle \cite{kaggle}, que simula o comportamento de clientes em um e-commerce. Este conjunto contém variáveis como tempo de navegação, frequência de compras, histórico de interações com a plataforma e outros dados importantes para modelar a probabilidade de um cliente realizar ou não uma compra em determinado período de tempo.

Portanto, este estudo visa explorar a aplicação da análise de sobrevivência via aprendizado de máquina, para prever comportamentos de compra em e-commerce, usando o Python como ferramenta de análise e o conjunto de dados mencionado para fornecer insights valiosos sobre os fatores que influenciam a decisão de compra dos consumidores online. A capacidade de prever esses comportamentos pode não só melhorar a experiência do cliente, mas também aumentar a eficiência das campanhas de marketing e da gestão de relacionamento com o cliente, resultando em benefícios para os negócios.




